\documentclass[11pt]{article}

\usepackage[margin=1in]{geometry}
\usepackage{amsmath,amssymb,amsthm}
\usepackage{enumitem}
\usepackage[hidelinks]{hyperref}
\urlstyle{same}

\newtheorem{theorem}{Theorem}
\newtheorem{proposition}[theorem]{Proposition}
\newtheorem{corollary}[theorem]{Corollary}
\newtheorem{lemma}[theorem]{Lemma}
\theoremstyle{definition}
\newtheorem{definition}[theorem]{Definition}
\theoremstyle{remark}
\newtheorem{remark}[theorem]{Remark}
\newtheorem{question}[theorem]{Question}

\newcommand{\F}{\mathbb{F}}
\newcommand{\Aff}{\operatorname{Aff}}
\newcommand{\Cut}{\operatorname{Cut}}
\newcommand{\Span}{\operatorname{span}}
\newcommand{\taglabel}[1]{\textnormal{\textbf{[#1]}}}
\newcommand{\PROVED}{\taglabel{PROVED}}
\newcommand{\TESTED}{\taglabel{TESTED}}
\newcommand{\STANDARD}{\taglabel{STANDARD MATH}}
\newcommand{\OPEN}{\taglabel{OPEN}}

\title{The FIFO linking theorem:\\affine response and the cut-space frontier}
\author{a9lim}
\date{August 2026}

\begin{document}
\maketitle

\begin{abstract}
This note attacks the general linking theorem underlying the
\textsc{echo}-\textsc{fifo}+dummy Gold-quadric realizer.  A finite graph is
played by opening vertices into a FIFO queue and later closing the front.  A
close flips a bit when the front has odd degree into the untouched set.  The
target theorem says that one isolated dummy lets the even player force even
total flip parity from either seat, for every graph.

The theorem is not proved here.  The contribution is an exact algebraic
reduction that replaces the failed local partner/debt induction by a global
affine-flow target.  Every play has a vector-valued score: its graph of disjoint
activity intervals.  For every fixed attacker strategy, the linking theorem is
equivalent to saying that zero lies in the $\F_2$-affine hull of all response
scores.  At a real FIFO front $f$, the edge space splits canonically as
\[
 E(K_L)=\Cut(K_L)\oplus E(K_{L\setminus\{f\}}),
\]
so openings before $f$ closes occupy one cut-space layer and all later play
occupies the smaller edge-space layer.  The quotient is the classical
two-graph coboundary chart: changing the front changes the affine-simplex
origin, and successive changes obey an absorption law.  This gives a canonical
filtration, but not a branchwise induction: an exact $P_4+d$ state has defender
children whose projected continuation affine hulls are disjoint.

Positive results accompany the reduction.  The vector live-star identity
proves the affine formulation, a parity-cell pairing theorem strictly extends
the earlier twin-pair solution, and the cut moment at the initial front
contracts affinely.  A second handshake refinement colours a vertex by its
degree parity and its parity of odd-degree neighbours: every colour class is
even on an even-order graph, and a same-colour pair balances all four joint
degree/incidence fibres.  A two-switch tail lemma closes the first genuine
adaptive-pairing obstruction.  Exact witnesses show why neither same-degree
pairing, greedy FIFO rotation, a bare symplectic quotient, nor a uniformly
bounded triangular fan closes the general case.  Scalar force sets give a
least-counterexample odd-routing lemma, while a stronger finite strategy has
an exact ranked repair-forest recursion: consecutive queue matchings may
extend, delete their front pair, or flip phase along an alternating path.
That flip carries a genuine two-graph curvature term.  Tetrahedral coherence
controls its label, but attacker pruning destroys the constant-coefficient
simplex contraction.  The remaining obligation is simultaneous cancellation
of the correlated cut and continuation moments by a strategy-relative chain.
The note therefore advances the proof
frontier without promoting the complete finite census through eight real
vertices to a proof.
\end{abstract}

\section{The reduced game}

Let $G=(R,E)$ be a finite simple graph on the \emph{real} vertices, and adjoin
one isolated vertex $d$.  A state is $(U,Q,\mathsf{ko})$, where $U$ is the set
of untouched vertices, $Q$ is an ordered FIFO queue of open vertices, and
$\mathsf{ko}$ records the one-move prohibition after a vertex is opened onto an
empty queue.  A legal move is
\begin{enumerate}[leftmargin=*,itemsep=1pt]
\item $\operatorname{OPEN}(x)$ for any $x\in U$, moving $x$ to the back of
      $Q$; or
\item $\operatorname{CLOSE}(f)$ for the front $f$ of $Q$, unless
      $\mathsf{ko}$ is set.
\end{enumerate}
A close of $f$ scores
\[
  \deg_U(f)\pmod 2.
\]
If no move is legal, the forced pass only clears $\mathsf{ko}$.  This can occur
only after $U$ is empty, when no later close can score.  The players have
opposite targets for the parity of the total score.

\PROVED{} Every vertex opens and closes exactly once.  FIFO makes the opening
and closing orders identical.  Consequently the activity intervals
$I_x=[o(x),c(x)]$ never properly nest.  For an edge $xy$ with $o(x)<o(y)$, the
edge contributes to the flip score precisely when $c(x)<o(y)$, i.e. precisely
when $I_x$ and $I_y$ are disjoint.  Thus the total score is
\begin{equation}\label{eq:scalar-score}
  F_G(h)=\sum_{xy\in E(G)}[I_x\cap I_y=\varnothing]\pmod2.
\end{equation}
The proper-interval vocabulary is standard; the adversarial parity problem is
not a standard queue-layout problem, because the queue stores vertices and the
play itself constructs their intervals.

\section{The vector score}

Write
\[
  \mathcal E_R=\F_2^{\binom{R}{2}}
\]
for the vector space whose coordinates are unordered pairs of real vertices.
For a completed history $h$, define its \emph{disjointness vector}
\[
  D(h)=\sum_{\{x,y\}\subseteq R}
       [I_x\cap I_y=\varnothing],xy\in\mathcal E_R.
\]
Pairing $D(h)$ with the adjacency vector $A_G$ recovers
\eqref{eq:scalar-score}.

At any moment let $L=U\cup Q$ be the live set.  For a real $x\in L$, put
\[
  s_L(x)=\sum_{y\in (L\cap R)\setminus\{x\}}xy,
\]
and set $s_L(d)=0$.  These are the vertex stars, or simplicial
coboundaries, of the live complete graph.

\begin{theorem}[vector live-star identity]\label{thm:live-star}
\PROVED{} Every completed history satisfies
\begin{equation}\label{eq:live-star}
  D(h)=\sum_{x\text{ opened}}s_{L_x}(x),
\end{equation}
where $L_x$ is the live set at the instant $x$ opens.
\end{theorem}

\begin{proof}
Fix a real pair $xy$ and suppose $x$ opens first.  The vertex $y$ is live when
$x$ opens, so $xy$ occurs once in $s_{L_x}(x)$.  If the activity intervals
overlap, $x$ is still live when $y$ opens and the coordinate occurs a second
time in $s_{L_y}(y)$; the two occurrences cancel over $\F_2$.  If the intervals
are disjoint, $x$ has already closed when $y$ opens and there is no second
occurrence.  The coordinate of the right-hand side is therefore exactly the
coordinate of $D(h)$.
\end{proof}

\begin{corollary}[scalar live-degree identity]
\PROVED{} For every graph $G$,
\[
  F_G(h)=\sum_{x\text{ opened}}\deg_{G[L_x]}(x)\pmod2.
\]
\end{corollary}

This upgrades the earlier scalar potential calculation to an identity before
a graph is chosen.  The graph enters only through the final linear functional
$z\mapsto\langle A_G,z\rangle$.

\section{Affine response is exactly the theorem}

Fix which seat seeks odd parity, and fix a deterministic strategy $S$ for that
attacker.  Move legality is graph-independent, so $S$ can be regarded as a
strategy in the schedule game before an adjacency vector is chosen.  Let
$\mathcal H_S$ be the set of terminal histories compatible with $S$ and with
arbitrary legal responses by the other player.

\begin{theorem}[affine separation]\label{thm:affine}
\PROVED{} The strategy $S$ forces odd score for some graph on $R$ if and only if
\[
  0\notin\Aff_{\F_2}\{D(h):h\in\mathcal H_S\}.
\]
Consequently, the general linking theorem is equivalent to the following
graph-free statement for both choices of attacker seat:
\begin{equation}\label{eq:affine-target}
  \text{for every }S,\qquad
  0\in\Aff_{\F_2}\{D(h):h\in\mathcal H_S\}.
\end{equation}
\end{theorem}

\begin{proof}
If $S$ forces odd score for a graph with adjacency vector $A$, then
$\langle A,D(h)\rangle=1$ for every $h\in\mathcal H_S$.  The whole affine hull
therefore lies in the affine hyperplane $\langle A,z\rangle=1$, which excludes
zero.

Conversely, write the affine hull as $a+W$.  If it excludes zero, then
$a\notin W$.  Elementary linear separation over $\F_2$ supplies a linear
functional $\ell$ with $\ell(W)=0$ and $\ell(a)=1$.  Every linear functional on
$\mathcal E_R$ is pairing with the adjacency vector of a simple graph on $R$.
For that graph, $\ell(D(h))=1$ for every response history, so $S$ forces odd
score.
\end{proof}

Equivalently, one must select an odd-cardinality formal multiset of response
histories whose disjointness vectors XOR to zero.  This is an
$\F_2$-flow problem: attacker nodes route all incoming weight along the move
prescribed by $S$, while defender nodes may split it among an odd number of
outgoing moves.  Uniformly summing all response leaves does \emph{not} work in
general; the flow must be selected.

There is a useful exact homogeneous recursion.  At a history $h$, let
$Z(h)\subseteq\F_2\oplus\mathcal E_R$ be the span of
$(1,D_h(\lambda))$ over response leaves $\lambda$ below $h$, where $D_h$ is
the future disjointness increment.  Give a close of a real front $f$ the
charge $c=\sum_{u\in U\cap R}fu$, and give an open, dummy close, or pass charge
zero.  This sums to $D(h)$ because a pair becomes disjoint exactly when its
earlier vertex closes while the later one is untouched.  Put
\[
  \tau_c(a,z)=(a,z+ac).
\]
Then a terminal node has $Z=\Span\{(1,0)\}$; an attacker node, unary after
fixing $S$, has $Z(h)=\tau_c Z(h')$; and a defender node has
\[
  Z(h)=\Span_m\tau_{c_m}Z(h_m).
\]
The desired affine response is exactly $(1,0)\in Z(\text{root})$.  Thus the
open problem is a label-preserving contraction of this recursion, not an
ordinary contraction of the response tree: ordinary boundaries have even
augmentation, while the required leaf chain has odd augmentation.

\TESTED{} The existing exhaustive graph census through eight real vertices
plus dummy, both attacker seats, implies \eqref{eq:affine-target} in those
dimensions by Theorem~\ref{thm:affine}.  Independent graph-free policy searches
also found no affine offset through eight real vertices; selected structured
policies at nine and ten real vertices likewise had zero offset.  These are
finite checks, not a proof of \eqref{eq:affine-target}.

\section{A solved parity-cell subclass}

The true/false-twin argument in the earlier draft can be weakened
substantially.

\begin{theorem}[parity-cell pairing]\label{thm:cells}
\PROVED{} Suppose all vertices, including the dummy if present, admit a perfect
partition $\mathcal M$ into two-element cells such that
\begin{equation}\label{eq:cell-condition}
  e_G(A,B)=0\pmod2
  \qquad\text{for every distinct }A,B\in\mathcal M.
\end{equation}
Then the second player forces even flip parity.
\end{theorem}

\begin{proof}
When the first player opens one member of a cell $A$, the second opens its
mate.  When the first player closes the FIFO-front member of $A$, its mate is
the next queue front and the second closes it.  Thus untouched vertices are
always a union of whole cells, and the queue is a concatenation of whole cells.
Ko causes no exception because the forced second opening completes the first
cell.

When the two members $a,a'$ of $A$ close consecutively, their two flip bits sum
to
\[
  \deg_U(a)+\deg_U(a')
  =\sum_{B\subseteq U}e_G(A,B)=0
\]
by \eqref{eq:cell-condition}.  Every flip is contained in one such consecutive
pair, so the total parity is even.
\end{proof}

Twin pairs satisfy \eqref{eq:cell-condition}, but the condition only asks for
even cross-parity between cells; no automorphism is required.

\TESTED{} This theorem is still not general.  Among the $1{,}044$
isomorphism classes on seven real vertices, $64$ boards obtained by adjoining
the isolated dummy have no such partition.  The first witness in the local
ordering has real edges
\[
  02,04,06,14,15,23.
\]
The general linking theorem nevertheless passes exact minimax on every one of
those boards.  Seidel switching cannot repair the pairing obstruction: between
two two-cells $A,B$, switching a set $S$ toggles
\[
 |A\cap S||B\setminus S|+|A\setminus S||B\cap S|
 \equiv ab+ab=0\pmod2.
\]

\section{FIFO gives a cut-space filtration}

Let $L\subseteq R$ be a live real set and let $f\in L$ be the FIFO front.  Let
$E(L)=\F_2^{\binom{L}{2}}$ and let
\[
  \Cut(L)=\Span\{s_L(x):x\in L\}.
\]
This is the cut space of the complete graph on $L$.

\begin{theorem}[front splitting]\label{thm:front-split}
\PROVED{} There is a canonical direct sum
\begin{equation}\label{eq:front-split}
  E(L)=\Cut(L)\oplus E(L\setminus\{f\}).
\end{equation}
The projection to the second summand is
\begin{equation}\label{eq:Tf}
  (T_fz)_{ij}=z_{ij}+z_{fi}+z_{fj},
  \qquad i,j\neq f.
\end{equation}
It kills every live star and is the identity on the canonical
$E(L\setminus\{f\})$ summand.
\end{theorem}

\begin{proof}
For $x=f$, the two incident coordinates in \eqref{eq:Tf} cancel; for
$x=i\neq f$, the $ij$ and $fi$ coordinates of $s_L(i)$ cancel.  Hence $T_f$
kills all generators of $\Cut(L)$.  If $z$ is supported away from $f$, then
$T_fz=z$.  Finally,
\[
 \dim\Cut(L)+\dim E(L\setminus\{f\})
 =(|L|-1)+\binom{|L|-1}{2}=\binom{|L|}{2},
\]
so the two subspaces are complementary and exhaust $E(L)$.
\end{proof}

There is a front-independent description of the same quotient.  Regard
$z\in E(L)$ as a simplicial $1$-cochain on the full simplex on $L$, and put
\[
  (\delta z)_{ijk}=z_{ij}+z_{ik}+z_{jk}.
\]
Then $\ker\delta=\Cut(L)$ and
\begin{equation}\label{eq:two-graph-chart}
  (T_fz)_{ij}=(\delta z)_{fij}.
\end{equation}
Thus $T_f$ is the chart based at $f$ of the odd-triple two-graph of $z$;
graphically it is Seidel complementation at $f$ followed by deletion of $f$
\cite{mallows-sloane,limouzy}.  Equivalently, if $B$ is the alternating form
with Gram matrix $z$ and
\[
  E_0=\{x\in\F_2^L:\textstyle\sum_i x_i=0\},\qquad
  b_i^{(f)}=e_i+e_f,
\]
then $T_fz$ is the Gram matrix of $B|_{E_0}$ in the basis
$\{b_i^{(f)}:i\ne f\}$.  Changing the basepoint from $f$ to $g$ gives
\[
 b_f^{(g)}=b_g^{(f)},\qquad
 b_i^{(g)}=b_i^{(f)}+b_g^{(f)}.
\]
In particular, after restricting away from distinct $f,g$,
\begin{equation}\label{eq:T-absorption}
  T_g(T_fz)=T_gz.
\end{equation}
This absorption law says that earlier quotient gauges cancel exactly; it does
not say that their affine payoff offsets cancel.

While $f$ remains front, opens do not change the live set.  By
Theorem~\ref{thm:live-star}, all their contributions lie in the current
$\Cut(L)$ layer.  Once FIFO deletes the real front $f$, every later contribution
is supported on $E(L\setminus\{f\})$.  Repeating the splitting along the forced
closure order gives a triangular decomposition of the complete play score.
This is the missing structural use of FIFO: the earlier scalar potential was
valid even if any queued vertex could close and therefore could not prove the
theorem.

The dummy needs a companion state in this induction.  Its star is zero, but
closing it does not shrink the real edge space.  A proof must use that zero
star to adjust affine augmentation and then transfer the odd response flow to
the unchanged real live set.  Treating the dummy as an ordinary deleted front
silently loses precisely the exceptional case.

\section{What the cone homotopy does and does not contract}

Fix a real front $f$ and a terminal continuation $\lambda$.  Let
$p_f(\lambda)\in\F_2^{L\setminus\{f\}}$ be the indicator of the real vertices
still untouched when $f$ closes, and let $W_f(\lambda)$ be the away-$f$
disjointness vector accumulated by the whole continuation.  Directly from the
definition of disjoint intervals,
\begin{equation}\label{eq:correlated-moment}
 D_h(\lambda)=\sum_i p_f(\lambda)_i\,fi+W_f(\lambda),
 \qquad
 T_fD_h(\lambda)=W_f(\lambda)+\delta p_f(\lambda),
\end{equation}
where $(\delta p)_{ij}=p_i+p_j$.  Hence for every formal response-leaf chain
$c$, of either augmentation,
\begin{equation}\label{eq:joint-cancellation}
 D_h(c)=0
 \quad\Longleftrightarrow\quad
 p_f(c)=0\ \text{ and }\ W_f(c)=0.
\end{equation}
The two-graph chart therefore replaces the desired edge moment by a
\emph{joint} cut/continuation moment; it does not discard the correlation.

One half of that joint target does contract cleanly at the root.

\begin{lemma}[initial cut-moment contraction]\label{lem:cut-contract}
\PROVED{} Suppose the attacker opens a real initial front $f$, so the isolated
dummy remains untouched and ko forbids the defender's immediate close.  For
every fixed attacker strategy up to the close of $f$,
\[
 0\in\Aff\{p_f(\lambda):\lambda\text{ is a compatible pre-close response}\}.
\]
\end{lemma}

\begin{proof}
If zero were absent, separation would give weights $a_i\in\F_2$ on the
remaining real vertices such that the attacker forced
$a\mathbin{\cdot}p_f=1$.  Give the dummy weight zero, and let $s$ be the total
weight of the still-unselected vertices.  On the ko-forced first response, the
defender selects the dummy when $s=0$, and selects a weight-one real vertex
when $s=1$; either choice leaves $s=0$.  On every later response, close $f$
when $s=0$, and select a weight-one vertex when $s=1$.  If the attacker selects
weight zero, $s$ remains zero and the defender closes.  If the attacker selects
weight one from a zero total, the number of remaining weight-one vertices was
positive and even, so another remains and the defender restores $s=0$.
Whenever the attacker closes $f$, therefore, $a\mathbin{\cdot}p_f=0$, a
contradiction.
\end{proof}

The scalar separation proof has a useful two-round refinement.  It repairs the
first explicit failure of degree-only pairing, although it still does not
iterate through an arbitrary FIFO prefix.

For a finite graph $H$, put
\[
 p_v=\deg_H(v)\pmod2,
 \qquad
 b_v=\sum_{u\sim_H v}p_u\pmod2,
 \qquad
 c(v)=(p_v,b_v).
\]
Thus $b_v$ is the parity of the odd-degree neighbours of $v$.

\begin{lemma}[two-bit handshake refinement]\label{lem:two-bit}
\PROVED{} If $H$ has even order, every one of the four colour classes
$c^{-1}(i,j)$ has even cardinality.  In particular every vertex has a
distinct same-colour mate.

If $c(v)=c(w)$ and $A=\{v,w\}$, put
$r_x=e_H(x,A)\pmod2$ for $x\notin A$.  Then each of the four sets
\[
 \{x\notin A:(p_x,r_x)=(i,j)\},\qquad (i,j)\in\F_2^2,
\]
also has even cardinality.
\end{lemma}

\begin{proof}
Let $O=\{v:p_v=1\}$ and $E=V(H)\setminus O$.  Handshaking gives $|O|=0$
in $\F_2$.  Moreover
\[
 \sum_{v\in O}b_v=2e_H(O)=0,
\]
so both $b$-fibres in $O$ are even.  The cut $e_H(O,E)$ is even, since
\[
 e_H(O,E)=\sum_{v\in O}\deg_H(v)=|O|=0
\]
in $\F_2$.  Hence the $b=1$ fibre in $E$ is even; $|E|$ is even because
$|V(H)|$ and $|O|$ are even, so the remaining fibre is even as well.

For the second assertion, the four parity moments of the map
$x\mapsto(p_x,r_x)$ on $V(H)\setminus A$ vanish:
\[
 \sum 1=0,
 \quad
 \sum p_x=0,
 \quad
 \sum r_x=p_v+p_w=0,
 \quad
 \sum p_xr_x=b_v+b_w=0.
\]
The internal edge $vw$, when present, contributes
$p_v+p_w=0$ to the last two displayed identities.  The four fibre
indicators are the four polynomials
$(1+p)(1+r),(1+p)r,p(1+r),pr$ over $\F_2$; their sums therefore vanish.
\end{proof}

\begin{remark}[what the second bit buys]
Suppose an attacker opens $v$ and the defender opens a same-colour mate $w$.
If the attacker closes $v$, closing $w$ makes the two flip charges cancel.
If the attacker instead opens $x$, Lemma~\ref{lem:two-bit} supplies a $y$
with both $p_y=p_x$ and $e_H(y,A)=e_H(x,A)$.  Thus the new cell matches the
opening-charge parity and is orthogonal to the first FIFO cell.  The equality
$b_v=b_w$ is exactly the missing joint moment in the first degree-only
counterexample.

This is a two-round statement, not an induction: the new cell can split the
four balanced fibres unevenly.  A later response may have to close rather
than open, which is the correlated continuation obstruction in
Question~\ref{q:contraction}.
\end{remark}

\TESTED{} Exact minimax on every graph isomorphism class through seven real
vertices plus the dummy gives a stronger root fact: whenever the same-colour
mate of Lemma~\ref{lem:two-bit} applies, \emph{every} such reply is winning
for the second-seat even player.  When the real graph has even order and the
attacker opens the unmatched dummy, every real second opening is winning
through six real vertices.  The first-seat even player also wins by opening
the dummy throughout the seven-vertex census.  The optional
\texttt{refinement} stage of \path{experiments/linking_game.py} reproduces
these checks.  They support the two-bit prefix but do not prove that its
balanced fibres survive later FIFO switches.

\section{Scalar response faces and the least counterexample}

It is useful to isolate exactly what an induction on a smaller even board
would have to prove.  Fix the player seeking odd parity and, for any game
state $h$, let $R(h)\subseteq\F_2$ be the set of future charge bits that this
attacker can force.  If a move of charge $c$ leads to $h'$, write
$c+R(h')=\{c+r:r\in R(h')\}$.  Backward induction gives
\begin{equation}\label{eq:force-recursion}
 R(h)=
 \begin{cases}
  \displaystyle\bigcup_m(c_m+R(h_m)),&h\text{ is an attacker node},\\[6pt]
  \displaystyle\bigcap_m(c_m+R(h_m)),&h\text{ is a defender node},
 \end{cases}
 \qquad R(\text{terminal})=\{0\}.
\end{equation}
Here a forced pass is the sole move from its state and has charge zero.
For a live vertex set $A$, use the following canonical roots:
\[
\begin{array}{c|c|c|c}
&Q&U&\text{mover / ko}\\ \hline
T(A)&()&A&\text{attacker / false}\\
K_x(A)&(x)&A-\{x\}&\text{defender / true}\\
S_x(A)&(x)&A-\{x\}&\text{defender / false}\\
Q_{xy}(A)&(x,y)&A-\{x,y\}&\text{attacker / false}.
\end{array}
\]
The same letters denote their force sets.  The boundary values are
$T(\varnothing)=K_x(\{x\})=\{0\}$.  Away from the empty-union boundary,
directly from
\eqref{eq:force-recursion},
\begin{equation}\label{eq:canonical-faces}
 T(A)=\bigcup_{x\in A}K_x(A)\quad(A\ne\varnothing),\qquad
 K_x(A)=\bigcap_{y\ne x}Q_{xy}(A)\quad(|A|\ge2),
\end{equation}
and
\begin{equation}\label{eq:S-face}
 S_x(A)=K_x(A)\cap
 \bigl(\deg_{A-\{x\}}(x)+T(A-\{x\})\bigr).
\end{equation}
Thus the scalar even-board statement is
\begin{equation}\label{eq:even-board}
 E(A):\qquad 1\notin T(A)\quad (|A|\text{ even}).
\end{equation}
It says that a first-seat odd attacker cannot force odd parity; by finite
determinacy, the second player can force even parity.

\begin{lemma}[least-counterexample odd routing]\label{lem:odd-routing}
\PROVED{} Suppose $A$ is a least-order even counterexample to
\eqref{eq:even-board}.  Then there is an $x\in A$ with $1\in K_x(A)$.  Put
\[
 Y=\{y\in A-\{x\}:\deg_A(y)=\deg_A(x)\pmod2\}.
\]
The set $Y$ has odd cardinality.  For every $y\in Y$, every target-$1$
strategy witnessing $1\in Q_{xy}(A)$ begins by opening a third vertex
$z_y$; it cannot begin by closing $x$.

At the resulting defender state $(x,y,z_y)$, such a strategy must satisfy
both
\begin{equation}\label{eq:routing-faces}
 1+\deg_{A-\{x,y,z_y\}}(x)\in Q_{y z_y}(A-\{x\})
\end{equation}
and target-$1$ forceability after every defender opening
$w\in A-\{x,y,z_y\}$.
\end{lemma}

\begin{proof}
Direct play proves $E(A)$ for $|A|=0,2$: on two vertices the forced opening
pair is followed by two zero-charge closes.  Hence a least counterexample has
order at least four, and Equation~\eqref{eq:canonical-faces} supplies $x$
without meeting either boundary case.  Each degree-parity class
in an even-order graph has even cardinality, so deleting $x$ leaves the
odd set $Y$.

Fix $y\in Y$.  If the attacker closes $x$ from $Q_{xy}(A)$, the charge is
\[
 c=\deg_{A-\{x,y\}}(x),
\]
and the child is $S_y(A-\{x\})$.  By minimality,
$T(A-\{x,y\})\subseteq\{0\}$, so \eqref{eq:S-face} gives
\[
 S_y(A-\{x\})
 \subseteq \{\deg_{A-\{x,y\}}(y)\}.
\]
But equality of the full degree parities gives
\[
 \deg_{A-\{x,y\}}(x)=\deg_{A-\{x,y\}}(y),
\]
where the possible edge $xy$ cancels from both sides.  The remaining target
after the close is $1+c$, not $c$, so this branch cannot force target $1$.
The attacker must therefore open some $z_y$.  The defender may now close
$x$ or open any remaining $w$.  Applying \eqref{eq:force-recursion} to
those children gives \eqref{eq:routing-faces} and the stated universal open
conditions.
\end{proof}

The odd routing is a real strengthening over the handshake lemma: it couples
an odd family of sibling histories.  It still cannot be collapsed one sibling
at a time.  A tempting twisted induction would label
$C=A-\{x,y\}$ by $a(u)=B(x+y,u)$ and ask for a $w$ with
\[
 1+a(z)+a(w)\notin Q_{zw}(C).
\]
\TESTED{} This is false first at $|C|=4$.  Take the path
$5-2-3-4$, put $(a_2,a_3,a_4,a_5)=(1,0,1,0)$, and $z=5$.
Exact recursion gives
\[
 T(C)=\{0\},\qquad
 Q_{52}=\{0\},\quad Q_{53}=\{0,1\},\quad Q_{54}=\{0\},
\]
while the three shifted targets are $0,1,0$.  All three survive.  This is
realized inside the six-vertex graph
\[
 E=\{03,12,13,14,23,25,34\},\qquad (x,y)=(0,1).
\]
Hence ordinary smaller-board induction, even augmented by one complete cut
bit, does not prove Lemma~\ref{lem:odd-routing}'s required joint sibling
cancellation.

The lemma cannot simply be followed by an independent quotient induction.
On real vertices $0,1,2$ plus $d$, fix a second-seat attacker so that the line
\[
 O_d,O_1,C_d,O_0,C_1,O_2,C_0,C_2
\]
is forced below the displayed defender choices.  At the $C_1$ branch from
$Q=(1,0),U=\{2\}$, one has $p_1=e_2$, $W_1=0$, $D=12$, and
$T_1D=02$.  The subtree below that branch has one leaf.  The fictitious $02$
terms in $W+\delta p$ cancel only after the cut and quotient pieces are
recombined; no compatible continuation chain realizes either term separately.
This is the smallest broken commuting diamond behind the false branchwise
induction.

\section{Why the deterministic local induction fails}

The live-degree pairing heuristic is correct only conditionally.  At a turn of
the block initiator the queue has even length.  If the initiator opens $x$, the
responder would like to open an untouched $y$ with the matching live-degree
parity.  If no such $y$ exists and the old queue is nonempty, closing the old
front rotates
\[
 (a_1,b_1,\ldots,a_r,b_r,x)
 \longmapsto
 (b_1,a_2,\ldots,a_r,b_r,x).
\]
Neither operation is a complete strategy.

\TESTED{} The smallest useful witnesses are exact minimax states, not
counterexamples to the linking theorem.
\begin{enumerate}[leftmargin=*,itemsep=3pt]
\item On $K_2+d$, after
\[
 A:O_d,\qquad D:O_0,\qquad A:C_d,
\]
the defender faces $Q=(0)$ and $U=\{1\}$.  The winning move is $O_1$; the
apparently natural mirror close $C_0$ loses.
\item On the real path $3-0-1-2$ plus $d$, after
\[
 A:O_0,\qquad D:O_1,\qquad A:O_d,
\]
the defender faces $Q=(0,1,d)$ and $U=\{2,3\}$.  Both available opens have the
wrong live-score parity, and closing $0$ loses.  The unique winning move is the
proactive wrong-parity open $O_3$, which creates an even odd-front corridor and
stores a debt before the local trap appears.
\end{enumerate}
Thus a state descriptor consisting only of current debt, front parity, parity
classes in $U$, and one FIFO rotation is insufficient.  The missing datum is a
recursive repair certificate in the untouched graph.  This is exactly what the
affine flow is allowed to encode globally.

\section{The exact local residue}

The scalar cut potential gives a complete local classification, provided its
limitation is stated explicitly.  Regard the queue as a set when counting
edges and put
\[
  P(Q,U)=e_G(Q,U)\pmod2.
\]
For a move $m$, let $\chi(m)$ be the change in $P$.  If
$R=Q\cup U$, direct edge accounting gives
\[
 \chi(O_x)=\deg_{G[R]}(x),\qquad
 \chi(C_f)=\deg_U(f),\qquad
 \chi(\mathrm{pass})=0.
\]
The middle term is the actual flip.

\begin{proposition}[charge-matching and the poison state]
\PROVED{} Suppose $P=0$, $|U|$ is even, and the position is nonterminal.
After an opponent open, there is always a distinct legal open with the same
$\chi$.  If a reply is due after an opponent close of charge $c$, there is a
legal response of charge $c$ except in exactly the following case:
\begin{equation}\label{eq:poison}
 \begin{aligned}
 c&=0,\qquad U\ne\varnothing,\qquad
 \deg_{R'}(u)=1\pmod2\quad\text{for every }u\in U,\\
 &Q'=\varnothing\quad\text{or}\quad
 \deg_U(\operatorname{front}Q')=1\pmod2.
 \end{aligned}
\end{equation}
where $Q'=Q\setminus\{f\}$ and $R'=Q'\cup U$ after the close.
\end{proposition}

\begin{proof}
At $P=0$,
\[
 \sum_{u\in U}\deg_R(u)=e(U,Q)=0\pmod2,
\]
because edges internal to $U$ are counted twice.  Hence the odd-charge opens
form an even set; since $|U|$ is even, the even-charge opens do too.  The charge
class containing the opponent's opener therefore contains another vertex.
Opening never has a ko obstruction, and the live set $R$ is unchanged by an
open, so this is a legal matching response.

If the opponent closes $f$ with charge $c$, then
\[
 \sum_{u\in U}\deg_{R'}(u)=e(U,Q')=c\pmod2.
\]
For $c=1$ an odd-charge open exists.  For $c=0$ an even-charge open exists
unless every untouched vertex has odd charge; in that exceptional case a
matching response still exists precisely when the new FIFO front exists and
has even degree into $U$.  A close clears ko, so that front close is legal.
This is exactly the complement of \eqref{eq:poison}.  When $U$ is empty all
remaining closes, including the close after the unique terminal pass, have
charge zero.
\end{proof}

Matching $\chi$ restores $P=0$, but does \emph{not} by itself pair the flip
score: a mixed odd-close/odd-open round contains one actual flip.  This is the
debt which the global chain must transport.  The dummy prevents
\eqref{eq:poison} while it is untouched (it is an even-charge open) or is the
new front (it is an even close), but it may be buried in the queue or already
closed.  The proposition therefore isolates the local obstruction without
pretending to solve its recurrence.

One portion of that recurrence is nevertheless well founded.  Let $g$ be the
accumulated flip debt and keep $P=e_G(Q,U)$ as above.

\begin{lemma}[fixed-front poison descent]\label{lem:poison-descent}
\PROVED{} Suppose it is the defender's turn, $g=1$, and the front $f$ has
positive even degree into $U$; ko may be set or clear.  The defender can open
a neighbour of $f$ so that either the debt is discharged within the next
attacker--defender round or the degree drops by exactly two.  If the result is
still positive the lemma may be repeated; if it is zero, the descended
anticomplete front can be closed but need not discharge the debt.  Hence the
fixed-front loop terminates.
Moreover
\[
 H=g+P
\]
is unchanged by every close.  Thus $g=1,H=0$ excludes a singleton
anticomplete front, although $H$ is not generally preserved by opens.
\end{lemma}

\begin{proof}
Opening one neighbour changes $\deg_U(f)$ from even to odd and clears ko,
since the queue was already nonempty.  If the attacker
closes $f$, that odd close discharges the debt.  If the attacker opens a
nonneighbour, the front remains odd and the defender closes it, again
discharging the debt.  The only remaining attacker move is to open a second
neighbour; then the front is even again, the defender faces the same local
problem when its degree remains positive, and exactly two of its untouched
neighbours have been removed.  At degree zero the neighbour-open loop stops;
ko is clear at this descended boundary, so the even close is legal and carries
no flip.

A close of charge $c$ changes both $g$ and $P$ by $c$, so it fixes $H$.
If $g=1$ and $H=0$, then $P=1$, whereas a singleton queue whose front is
anticomplete to $U$ has $P=0$.  Finally an open $x$ changes $H$ by
$\deg_{G[Q\cup U]}(x)$, which may be odd; therefore this exclusion is a local
corridor condition, not a global invariant.
\end{proof}

The first obstruction to same-action pairing does have a complete local
repair.  Write $B(u,v)$ for adjacency in the current graph.

\begin{lemma}[two-switch tail]\label{lem:two-switch-tail}
\PROVED{} Suppose the defender moves just after an attacker close, the score
debt is zero, and
\[
 U=\{x,y\},\qquad Q=(v_1,\ldots,v_{2r+1}).
\]
If
\begin{equation}\label{eq:tail-pairs}
 B(v_{2i-1},y)=B(v_{2i},y)\quad(1\le i\le r),
 \qquad B(v_{2r+1},y)=B(x,y),
\end{equation}
then the defender forces even final score with at most two changes of action
type relative to the attacker's preceding move.
\end{lemma}

\begin{proof}
The defender first opens $x$, so the even queue is paired as
$(v_1,v_2),\ldots,(v_{2r+1},x)$ and every pair has equal adjacency to the sole
untouched vertex $y$.  While the attacker closes, the defender closes the next
front; the two flip bits in each pair cancel by \eqref{eq:tail-pairs}.  If the
attacker opens $y$, the defender closes the front.  Now $U$ is empty and every
later close has charge zero.  If the attacker never opens $y$ before the pairs
are exhausted, $y$ is the harmless terminal singleton.  The initial
close/open and the possible later open/close are the two switches.
\end{proof}

\TESTED{} Same-action adaptive pairing first fails, without a dummy, on two
of the seven exceptional order-eight graphs: \verb|GCZN^{| and
\verb|GEjt~{|.  Call them $G_1,G_2$ in the table below.
Exact minimax gives minimum worst-case switch count two on each.  Their worst
paths enter Lemma~\ref{lem:two-switch-tail} at
\[
\begin{array}{c|c|c|c}
&Q&U&(x,y)\\ \hline
G_1&(1,2,4,3,5)&\{6,7\}&(6,7)\\
G_2&(1,2,3,4,5)&\{6,7\}&(6,7).
\end{array}
\]
In both graphs $7$ dominates the remaining live graph, so
\eqref{eq:tail-pairs} is immediate.  This proves the tail repair, not the
still-missing prefix theorem that would force such a tail on every board.

\section{The zero-normalized kernel and dynamic queue matchings}

The order-eight search reveals a stronger scalar target than
\eqref{eq:even-board}.  Sample play after every defender move, with the
odd-seeking player about to move and accumulated flip score zero.  Away from
the unique terminal-pass boundary, the queue then has even length.  Put
\[
 \rho(U,Q)=|U|+\frac{|Q|}{2};
\]
every two-move block decreases $\rho$ by one.

\begin{proposition}[zero-normalized repair recursion]
\label{prop:repair-forest}
\PROVED{} Let $Q$ have even length.  There is a strategy returning the flip
score to zero after every defender move if and only if the following recursive
predicate $\mathsf{Safe}(U,Q)$ holds.
\begin{itemize}[leftmargin=*]
\item $\mathsf{Safe}(\varnothing,())$ holds.
\item $\mathsf{Safe}(\{z\},())$ holds as the terminal macro
      $O_z$, forced zero-charge pass, $C_z$.  The pass is the defender's move;
      after it the score is still zero, and the last close also has charge
      zero.
\item After an attacker open $O_z$, a legal zero-charge reply is either
      $O_w$ for some $w\in U-\{z\}$, leading to
      \[
       \mathsf{Safe}(U-\{z,w\},Qzw),
      \]
      or, when $Q=(f,Q_0)$ and
      $\deg_{U-\{z\}}(f)=0$, the close $C_f$, leading to
      \[
       \mathsf{Safe}(U-\{z\},Q_0z).
      \]
\item After an attacker close $C_f$ of charge
      $\epsilon=\deg_U(f)$ from $Q=(f,h,Q_0)$, a same-charge reply is
      $C_h$ when $\deg_U(h)=\epsilon$, leading to
      $\mathsf{Safe}(U,Q_0)$.  When $\epsilon=0$, the defender may instead
      open $w\in U$, leading to
      $\mathsf{Safe}(U-\{w\},hQ_0w)$.
\end{itemize}
Except inside the displayed terminal macro, every legal attacker move must
have at least one listed reply whose child is safe.
\end{proposition}

\begin{proof}
All opens and passes have flip charge zero, and the only nonzero move charge is
the displayed untouched degree of a FIFO front.  Consequently the list
contains exactly the replies whose charge equals the preceding attacker
charge.  The singleton macro is the only time the defender has no ordinary
reply: opening the last untouched vertex onto an empty queue creates ko, so
the forced pass clears it and leaves one final zero-charge close.  Every other
child has rank $\rho-1$.  Induction on $\rho$, treating the macro as a terminal
leaf, proves sufficiency by selecting the recorded child after every attacker
move.  Conversely, record the first reply of any strategy which restores zero
after each defender move, and apply the same argument to every reachable
child.  This produces the recursion.
\end{proof}

The resulting certificate is a ranked, branch-dependent repair forest.  It is
not yet a structural existence proof: defining the forest by the winning
strategy would be circular.  It does, however, state exactly what the finite
search certifies and prevents weaker local invariants from being mistaken for
the theorem.

\TESTED{} Every even-order graph through order eight has
$\mathsf{Safe}(V,())$.  On all $12{,}346$ order-eight isomorphism classes, the
second player can therefore keep the \emph{actual} flip score zero after every
one of her moves.  The first failures of the stronger same-action rule occur
on \verb|GCZN^{| and \verb|GEjt~{|; the zero-normalized strategy survives by
using a close/open or open/close rotation.  Static selectors fail more
sharply: on \verb|G?bFf[| an opener's unique same-$(p,b)$ mate is unsafe, and
on \verb|GCz^vw| an odd-degree opener's only safe reply has even degree.
These are failures of proposed selectors, not of
Proposition~\ref{prop:repair-forest} or the linking theorem.

The rotations have an exact ordered-matching interpretation.  At a checkpoint
write
\[
 Q=(v_0,v_1,\ldots,v_{2r-1})
\]
and match consecutive queue vertices.  Pair opens extend this matching and
pair closes delete its front edge.  If the attacker appends
$z=v_{2r}$ and the defender gives the legal zero close $C_{v_0}$, the new
matching is shifted by one place.  Let
\[
 P(W)=\sum_{i=0}^{2r-1}[v_i,v_{i+1}],
 \qquad W=(v_0,\ldots,v_{2r}).
\]
Then the old and shifted matching chains differ by $P(W)$ and
\[
 \partial P(W)=e_{v_0}+e_z.
\]
More precisely, with the triangle fan
\[
 F(W)=\sum_{i=1}^{2r-1}[v_0,v_i,v_{i+1}],
\]
one has
\begin{equation}\label{eq:phase-fan}
 P(W)+[v_0,z]=\partial F(W).
\end{equation}
For the adjacency $1$-cochain $B$, therefore,
\begin{equation}\label{eq:phase-holonomy}
 \langle B,P(W)\rangle
 =B(v_0,z)+
 \sum_{i=1}^{2r-1}\delta B(v_0,v_i,v_{i+1}).
\end{equation}
The second term is genuine two-graph holonomy.  Endpoint exchange, a queue
matching span, or a symplectic basis invariant sees only $B(v_0,z)$ and is
therefore insufficient on an arbitrary graph.

The triangle shears are coherent but not flat.  If
$\tau(i,j,k)=\delta B(i,j,k)$, then
\begin{equation}\label{eq:tetrahedron-coherence}
 \tau(j,k,l)+\tau(i,k,l)+\tau(i,j,l)+\tau(i,j,k)=0,
\end{equation}
which is $\delta^2B=0$.  Thus the ambient complete simplex fills the fan in
\eqref{eq:phase-fan} consistently.  The adversarial obstruction is in its
coefficients: after fixing an attacker strategy, compatible histories are not
closed under deleting an interior letter.  Defender coefficients may split,
but an attacker coefficient must follow one prescribed child.  Projection to
compatible histories consequently fails to commute with the simplicial
boundary.

This is not repaired by ordinary topology of injective words.  Those words
form a Boolean cell complex with
\[
 \partial[v_0\cdots v_k]
 =\sum_i[v_0\cdots\widehat v_i\cdots v_k]
 \quad\text{over }\F_2
\]
\cite{jonsson-welker,gan}.  Define the strategy carrier here as the Boolean
subcomplex generated by the attacker-compatible pre-close tail words.  Even
with a dummy, attacker pruning need not leave this carrier acyclic.  On
$K_2+d$, let the attacker open a real front $f$ and, after the forced-open
reply is chosen as $O_a$ or $O_d$, close $f$ immediately.  Before the first
close the carrier consists of the two isolated vertices $[a]$ and $[d]$; the
joining cells $[ad]$ and $[da]$ have been pruned.  Its reduced
$H_0$ is nonzero.  The linking response still exists because the two close
charges differ.  Hence any successful carrier theorem must retain the charge
and continuation labels; ordinary shellability or a constant-coefficient cone
is not enough.

Three further exact obstructions delimit the remaining Boolean contraction.
First, the two-bit colour is a prefix moment, not a winning selector.
\TESTED{} On the eight-vertex graph
\[
\begin{split}
 E=\{&01,03,05,07,12,14,15,17,24,26,27,34,\\
     &37,45,46,47,57,67\},
\end{split}
\]
vertices $6,7$ have the same $(p,b)$ colour, but $Q_{67}$ is unsafe while
$Q_{76}$ is safe.  The unique first spoiler in $Q_{67}$ is $O_5$.
There is also an opener with no safe same-colour mate: in graph6 class
\verb|G?ben[|, with
\[
 E=\{04,05,15,25,06,16,36,56,07,17,27,47,57,67\},
\]
vertex $2$ has sole same-colour mate $7$, yet $Q_{27}$ is unsafe.  Its
unique first spoiler is $O_3$; the reverse $Q_{72}$ is safe.  Thus neither
orientation nor existential choice within one colour class proves the root
theorem.

Second, a safe residual strategy cannot always be lifted by insisting on an
open response.  On
\[
\begin{split}
 E=\{&01,02,03,04,05,06,12,13,14,17,23,26,\\
     &27,35,37,56,57\},
\end{split}
\]
the checkpoint
\[
 Q=(4,3),\qquad U=\{0,1,2,5,6,7\}
\]
is safe in the zero-normalized recursion.  After attacker $O_6$, the unique
safe reply is the zero close $C_4$, leading to
$Q=(3,6)$ and $U=\{0,1,2,5,7\}$.  The five open replies by
$0,1,2,5,7$ are all unsafe; respective short spoilers may be chosen as
$C_4,C_4,O_7,O_7,O_1$.  Hence the phase-changing OC branch is not an
artifact of the finite search: it is indispensable even when $|U|$ was even.

Finally, time reversal does not repair the strategic pruning.  For a
pass-free complete event word put
\[
 \rho(e_1\cdots e_m)=\bar e_m\cdots\bar e_1,
 \qquad \overline{O_v}=C_v,\quad \overline{C_v}=O_v.
\]
\PROVED{} This is a legal FIFO schedule, swaps the two seats, and preserves
$D$ coordinatewise.  Indeed, if both opening and closing order in the original
word are $v_1,\ldots,v_n$, both orders after reversal are
$v_n,\ldots,v_1$, so FIFO is preserved.  A pass-free word has no adjacent
singleton block $O_vC_v$, hence reversal creates no ko violation.  Position
$i$ moves to $2n+1-i$, which swaps seat parity, while each activity interval
$[o_v,c_v]$ reflects to
$[2n+1-c_v,2n+1-o_v]$; interval disjointness, and therefore $D$, is
unchanged.

This identity does not extend to a strategy contraction.  If one attempts to
set $\rho(P)=P$, a terminal singleton component $O_x,P,C_x$ moves to the
beginning; when other vertices remain untouched there, its pass is not forced
and is illegal.  Even without passes, fixed strategies are anti-causal under
$\rho$.  Put the original attacker in the even positions.  One deterministic
attacker strategy contains both histories below: after $O_0$ it prescribes
$O_d$; after the defender's choice $C_0$ or $O_1$ it prescribes respectively
$O_1$ or $C_0$; and after $C_d$ it prescribes $C_1$.  The smallest such local
witness has real vertices $0,1$ plus $d$:
\[
\begin{array}{c|c}
h_C=O_0O_dC_0O_1C_dC_1&D(h_C)=01,\\
h_O=O_0O_dO_1C_0C_dC_1&D(h_O)=0.
\end{array}
\]
The two histories reconverge across the commuting $C_0/O_1$ diamond, whose
label difference is the edge coordinate $01$.  Reversal gives a common
prefix $O_1O_d$ but demands different moves, $C_1$ and $O_0$, at the same
odd-position attacker node.  Hence the reversed pair cannot lie under one
deterministic reversed attacker strategy.  Separately, the diamond involution
$\iota$ exchanging $C_0$ and $O_1$ swaps $h_C$ and $h_O$, so
\[
 \F_2\{h_C,h_O\}^{\iota}=\{0,h_C+h_O\}.
\]
Both invariant chains have even augmentation (and
$D(h_C+h_O)=01$), whereas the valid odd zero-moment certificate is the
nonsymmetric singleton $h_O$.  Diamond symmetrization therefore destroys the
odd augmentation rather than proving it.

\section{The remaining global lemma}

The proof target can now be stated without Gold forms, graph enumeration, or a
particular local menu.

\begin{question}[global FIFO affine contraction]\label{q:contraction}
\OPEN{} Fix either attacker seat and a deterministic attacker strategy $S$ in
the schedule game with one isolated dummy.  Does there always exist an odd
$\F_2$-flow through the defender branches whose terminal disjointness moment is
zero, with the flow allowed to couple different defender children and different
levels of the front filtration?
\end{question}

An affirmative answer proves the general linking theorem by
Theorems~\ref{thm:affine} and~\ref{thm:front-split}.  The tempting stronger
claim---that projected continuation hulls of different defender children meet
one common coset---is false.  At the $P_4+d$ state above, fix a
history-dependent attacker which follows different continuations after $O_2$
and $O_3$.  Under the front projection $T_0$ to the edge coordinates
$12,13,23$, exact recursion gives the two open-child continuation hulls
\[
  \{0\},
  \qquad
  \Aff\{23,12+23\}.
\]
They are disjoint.  Exhaustive recursion over the smaller states with two real
open children shows that this is the first such childwise failure.

Thus a valid contraction must transport nonzero target cosets and cancel them
\emph{across} defender children; it cannot solve every projected child as an
independent zero-target subproblem.  In particular, an induction classified
only by dummy presence, live-set parity, attacker seat, and front position is
false on reachable states.  The local relations
\[
  s_L(d)=0,
  \qquad
  \sum_{x\in L\cap R}s_L(x)=0
\]
provide the cone and the cut dependence, but not the required global coupling.

This boundary also explains why the nearest imported techniques do not close
the problem.  A standard strategy-stealing pass must be optional and harmless;
the dummy has two mandatory events which change FIFO legality.  A static
copycat argument requires an involution or a parity-cell partition, absent on
the $64$ finite witnesses above.  Local complementation changes the required
pairing data but is not a legal FIFO move.  Queue layouts supply terminology
for nonnesting but no adversarial parity strategy.

\TESTED{} A useful stronger conjecture is false: without the dummy, the
initial mover does not always force even score on even-order graphs.  The first
failures occur at order eight, in exactly seven isomorphism classes:
\verb|GCZN^{|, \verb|GCrU~{|, \verb|GEhvn{|, \verb|GCx}~{|,
\verb|GEjt~{|, \verb|GEl~v{|, and \verb|GUZv^{|.  They are
anti-mover-controlled: the second player can force the target parity, whether
that target is even or odd.  In particular the weaker finite conjecture that a
second-seat even player wins every even-order no-dummy board survives this
screen, but it is not proved.  Bare alternating rank or Witt class cannot
classify these outcomes: $P_3$ and $K_2\sqcup K_1$ have congruent rank-two
alternating adjacency forms over $\F_2$ but different seat outcomes.

\PROVED{} The smallest genuine affine cancellation is already a three-history
fan, not a dummy cone.  With real vertices $1,2,3$, let a first-seat attacker open the
dummy and thereafter close whenever possible, otherwise opening the least
untouched vertex.  After the shared prefix $O_d,O_1,C_d$, the defender can
realize the three compatible terminal histories
\[
\begin{array}{c|l|c}
 &\text{remaining moves}&D(h)\\ \hline
(i)&O_2,C_1,O_3,C_2,C_3&\{13\}\\
(ii)&O_3,C_1,O_2,C_3,C_2&\{12\}\\
(iii)&C_1,O_2,O_3,C_2,C_3&\{12,13\}.
\end{array}
\]
Their three vectors XOR to zero, although none is zero.  In (iii), $O_2$
opens an empty queue and $O_3$ is the ko-forced response, so all three rows are
legal.  The cancellation happens after the dummy has closed.

\TESTED{} Three histories do not suffice uniformly.  With five real vertices
plus $d$, let a first-seat attacker close whenever legal and otherwise open the
least untouched vertex.  Its $132$ distinct response scores contain neither
zero nor three vectors with XOR zero.  The minimum odd certificate has support
five; one is
\[
\begin{aligned}
&\{01,02,03,13,14,24\},
&&\{01,02,04,12,13,23,24\},\\
&\{01,03,04,12,14,23\},
&&\{02,03,04,12,14,23\},\\
&\{01,02,03,04,12,14,23\}.
\end{aligned}
\]
The five rows XOR to zero.  Thus the triangular fan is a real local
cancellation, not a bounded global normal form for affine response flows.

\TESTED{} Even the natural fan over \emph{all} second openings is too rigid.  On real
vertices $0,1,2$ plus $d$, let the attacker first play $O_0$; after replies
$O_1$ or $O_2$ use ``open least untouched, otherwise close'', and after $O_d$
use ``close whenever possible, otherwise open least''.  Exact terminal affine
hulls conditioned on the second reply, in coordinates $01,02,12$, are
\[
  \mathcal A_1=\{0\},\qquad
  \mathcal A_2=\{0\},\qquad
  \mathcal A_d=\{3,7\}.
\]
The immediate star charges of the three branches XOR to zero, but no
one-leaf-per-branch transversal does.  The global affine target survives by
retaining only an odd subflow inside $\mathcal A_1$ or $\mathcal A_2$.  Hence a
proof must choose its odd subset of defender branches using continuation data;
neither dummy-edge contraction nor the full root fan is a valid fixed recipe.

\section{Consequence for the Gold realizer}

\PROVED{} The reductions in \path{writeups/goldarf.tex} identify the terminal
$\sigma$-charge with the overlap parity of $G$-edges.  Since overlap and
disjointness partition $E(G)$,
\[
  \sigma(h)=|E(G)|+F_G(h)\pmod2.
\]
Therefore an affirmative answer to Question~\ref{q:contraction} would force
$\sigma=|E(G)|$ for every support graph and both seats.  On a Gold support
graph this is the required quadratic value, upgrading the verified
$m\in\{4,8\}$ realizer to every $m$.

\OPEN{} This would still solve only the mechanism half of the original
\texttt{tis} problem.  The realizer is $\sigma$-valued; recasting its forced
charge as a natural normal-, mis\`ere-, or loopy-outcome P-set remains a
separate open step.

\section{Verification boundary}

The following claims have different status and should not be flattened.
\begin{itemize}[leftmargin=*]
\item \PROVED{} The interval/disjointness reduction, vector live-star identity,
      affine separation equivalence, homogeneous response recursion,
      parity-cell theorem, front cut-space splitting and two-graph absorption,
      correlated-moment identity, initial cut contraction, charge/poison
      classification, two-switch tail, scalar force-set recursion,
      least-counterexample odd routing, zero-normalized repair recursion,
      queue-matching phase-fan identity, tetrahedral coherence, and displayed
      triangular fan and fixed-front poison descent, together with the
      pass-free reversal identity and its local causal obstruction, are
      complete proofs in this note.
\item \TESTED{} Exact minimax verifies the general linking theorem on every
      isomorphism class through eight real vertices plus dummy, both seats, and
      verifies the two local-strategy witnesses.  The parity-cell census is
      complete through seven real vertices.
\item \TESTED{} The childwise common-coset strengthening is false, first at the
      displayed $P_4+d$ state.  This is a failure of the proposed induction,
      not of the linking theorem.
\item \TESTED{} The mandatory full second-opening fan also fails at three real
      vertices plus dummy, despite the immediate star identity.  The global
      affine target survives inside a single reply branch.
\item \TESTED{} Same-action adaptive pairing first fails on two order-eight
      no-dummy boards and needs exactly two switches there.  A three-leaf
      affine certificate first fails for the displayed five-real-vertex
      strategy, whose exact minimum support is five.  These are failures of
      stronger proof templates, not of the linking theorem.
\item \TESTED{} Every even-order graph through order eight has a ranked
      zero-normalized repair forest.  Static degree, two-bit, matching,
      symplectic-flag, and one-cut selectors have the explicit failures stated
      above.  The finite forests do not prove their own existence at arbitrary
      order.
\item \TESTED{} Ordered same-colour safety and the weaker existence of a safe
      same-colour mate both fail at order eight.  The open-completion lift also
      fails at the displayed safe checkpoint, where a zero close is the unique
      safe response.  These are failures of structural proof templates, not
      counterexamples to the root theorem.
\item \OPEN{} The global FIFO affine-contraction statement
      \eqref{eq:affine-target} is not proved.  Neither the complete finite
      census nor the graph-free affine searches are evidence of a general
      chain contraction.
\end{itemize}

\begin{thebibliography}{9}
\bibitem{heath-rosenberg}
L.~S. Heath and A.~L. Rosenberg,
\emph{Laying out graphs using queues},
SIAM Journal on Computing 21 (1992), 927--958,
\href{https://doi.org/10.1137/0221055}{doi:10.1137/0221055}.

\bibitem{proper-interval}
P.~Klav\'ik, Y.~Otachi, and J.~\v{S}ejnoha,
\emph{On the classes of interval graphs of limited nesting and count of
lengths},
Algorithmica 81 (2019), 1490--1511,
\href{https://arxiv.org/abs/1510.03998}{arXiv:1510.03998}.

\bibitem{strategy-stealing}
G.~Bodwin and O.~Grossman,
\emph{Strategy-stealing is non-constructive},
\href{https://arxiv.org/abs/1911.06907}{arXiv:1911.06907}.

\bibitem{mallows-sloane}
C.~L. Mallows and N.~J.~A. Sloane,
\emph{Two-graphs, switching classes and Euler graphs are equal in number},
SIAM Journal on Applied Mathematics 28 (1975), 876--880,
\href{https://doi.org/10.1137/0128070}{doi:10.1137/0128070}.

\bibitem{limouzy}
V.~Limouzy,
\emph{Seidel minor, permutation graphs and combinatorial properties},
\href{https://arxiv.org/abs/0904.1923}{arXiv:0904.1923}.

\bibitem{jonsson-welker}
J.~Jonsson and V.~Welker,
\emph{Complexes of injective words and their commutation classes},
\href{https://arxiv.org/abs/0712.2143}{arXiv:0712.2143}.

\bibitem{gan}
W.~L.~Gan,
\emph{Complex of injective words revisited},
\href{https://arxiv.org/abs/1608.04496}{arXiv:1608.04496}.
\end{thebibliography}

\end{document}
